\begin{table}[h!]
\centering
\small
\caption{Prospective ORP sizing examples under independent-unit assumptions.}
\label{tab:app-orp-prospective-sizing}
\begin{minipage}[t]{0.47\textwidth}
\centering
\textbf{Panel A: equal subgroup sizes}\par\vspace{2pt}
\scriptsize
\begin{tabular}{@{}rrr@{}}
\toprule
\textbf{Absolute gap $\delta$} & \textbf{$n$/group, 80\%} & \textbf{$n$/group, 90\%} \\
\midrule
0.05 & 903 & 1208 \\
0.10 & 248 & 332 \\
0.15 & 119 & 159 \\
0.20 & 71 & 95 \\
\bottomrule
\end{tabular}
\end{minipage}
\hfill
\begin{minipage}[t]{0.47\textwidth}
\centering
\textbf{Panel B: detection of at least one failure}\par\vspace{2pt}
\scriptsize
\begin{tabular}{@{}rrrr@{}}
\toprule
\textbf{Failure rate $q$} & \textbf{$n$, 80\%} & \textbf{$n$, 90\%} & \textbf{$n$, 95\%} \\
\midrule
0.1\% & 1609 & 2302 & 2995 \\
1.0\% & 161 & 230 & 299 \\
5.0\% & 32 & 45 & 59 \\
\bottomrule
\end{tabular}
\end{minipage}

\vspace{4pt}
\parbox{0.98\textwidth}{\footnotesize
\textit{Notes:} Panel A assumes $p_{\mathrm{light}}=0.85$,
$p_{\mathrm{dark}}=0.85-\delta$, a two-sided $\alpha=0.05$ test, equal group
sizes, and the large-sample two-proportion approximation. Panel B is exact for
a zero-acceptance plan: $n$ is the number of independent eligible units needed
for probability $1-\beta$ of observing at least one failure when the failure
rate is $q$. Clustering, unequal weights, or non-exchangeable curation require
larger nominal samples or simulation under the intended design.}
\end{table}
